% English draft based on appendices/appendix_e_emergence_homology_ja.tex
\section{Emergence Homology}

The emergence schema used in this paper can be written:
\[
U=\Phi(L,H_{ij}).
\]
Here $L$ is lower-layer motion, $H_{ij}$ is mutual constraint, and $U$ is an upper-layer attractor. The upper layer does not erase the lower layer; it stabilizes its motion as a new operating level.

\begin{longtable}{@{}p{0.20\linewidth}p{0.30\linewidth}p{0.35\linewidth}@{}}
\toprule
Domain & Lower motion $L$ & Upper attractor $U$ \\
\midrule
\endfirsthead
\toprule
Domain & Lower motion $L$ & Upper attractor $U$ \\
\midrule
\endhead
Physics & causal logs and closure dynamics & particle-like attractor \\
Life & informational flow and state development & $\Omega_{\text{life}}$ \\
Intelligence & recursive log motion $L_R$ & emergent stopping layer $L_E$ \\
Society & interaction among individual $L_E$ layers & $\Phi_{\text{net}}$ \\
\bottomrule
\end{longtable}

The same emergence grammar appears at multiple layers. This does not collapse the layers into one another. It shows that VED uses a repeated pattern of non-closure, constraint, and attractor formation.

\bigskip
